\documentclass[12pt,aspectratio=1610]{ecnubeamerformyself}

% Override image paths (set here so the .cls keeps defaults)
% \renewcommand{\ECNUBackgroundImage}{img/bg.pdf}
% \renewcommand{\ECNUFrameLogoImage}{img/logo-white.pdf}
% \renewcommand{\ECNUFirstPageBG}{img/firstpagebg.jpg}
% \renewcommand{\ECNUMathLogo}{img/math-logo.pdf}
% \renewcommand{\ECNULibThanks}{img/lib.png}

\addbibresource{Accessories/ref.bib}
% Canonical notation definitions with and without unicode-math.
% Input this file in the preamble, after loading unicode-math when it is used.
\makeatletter
\@ifpackageloaded{unicode-math}{%
	\newcommand{\notationmathbold}[1]{\symbf{#1}}%
	\newcommand{\notationmathroman}[1]{\symrm{#1}}%
	\newcommand{\notationmathupright}[1]{\symup{#1}}%
	\newcommand{\notationmathbb}[1]{\mathbb{#1}}%
	\newcommand{\notationmathcal}[1]{\mathcal{#1}}%
	\newcommand{\notationmathscr}[1]{\mathscr{#1}}%
	\newcommand{\notationmathcaloperator}[1]{\mathcal{#1}}%
}{%
	\RequirePackage{mathrsfs}%
	\RequirePackage{bbm}%
	\newcommand{\notationmathbold}[1]{\mathbf{#1}}%
	\newcommand{\notationmathroman}[1]{\mathrm{#1}}%
	\newcommand{\notationmathupright}[1]{\mathrm{#1}}%
	\newcommand{\notationmathbb}[1]{\mathbbm{#1}}%
	\newcommand{\notationmathcal}[1]{\mathit{#1}}%
	\newcommand{\notationmathscr}[1]{\mathit{#1}}%
	\newcommand{\notationmathcaloperator}[1]{%
		\expandafter\notationmathcaloperator@classic#1\@nil
	}%
	\def\notationmathcaloperator@classic#1#2\@nil{%
		\mathcal{#1}\mathit{#2}%
	}%
}
\makeatother

% Customize operator font
\newcommand{\upoperatorname}[1]{\mathop{\notationmathupright{#1}}\nolimits}
\newcommand{\caloperatorname}[1]{\mathop{\notationmathcaloperator{#1}}\nolimits}
\newcommand{\frakoperatorname}[1]{\mathop{\mathfrak{#1}}\nolimits}
% \newcommand{\bfoperatorname}[1]{\mathop{\notationmathbold{#1}}\nolimits}
\newcommand{\uplimoperatorname}[1]{\mathop{\notationmathupright{#1}}\displaylimits}
% limits, nolimits, and displaylimits are used to control the position of subscript and superscript



% general functions and symbols
\newcommand{\kk}{{\notationmathbold{k}}}
\newcommand{\kkk}{{\notationmathbb{k}}}
\newcommand{\KK}{{\notationmathbold{K}}}
\newcommand{\KKK}{{\mathbb{K}}}
\newcommand{\rkk}{{\kappa}} % residue field
\newcommand{\fkk}{{\mathscr{K}}} % fraction field

\newcommand{\im}{\sqrt{-1}} % imaginary unit
\newcommand{\Real}{\upoperatorname{Re}}

\newcommand{\ten}{\otimes}

\newcommand{\injmap}{\hookrightarrow}
\newcommand{\surjmap}{\twoheadrightarrow}
\newcommand{\congmap}{\xrightarrow{\sim}}
\newcommand{\isomap}{\xrightarrow{\sim}}

\newcommand{\invlim}{\varprojlim}
\newcommand{\dirlim}{\varinjlim}

\renewcommand{\lim}{\uplimoperatorname{lim}}
\newcommand{\colim}{\uplimoperatorname{colim}}
\renewcommand{\limsup}{\uplimoperatorname{lim\,sup}}
\renewcommand{\liminf}{\uplimoperatorname{lim\,inf}}
\renewcommand{\max}{\uplimoperatorname{max}}
\renewcommand{\min}{\uplimoperatorname{min}}
\renewcommand{\sup}{\uplimoperatorname{sup}}
\renewcommand{\inf}{\uplimoperatorname{inf}}

\newcommand{\id}{{\upoperatorname{id}}}

\newcommand{\vol}{\upoperatorname{vol}}
\newcommand{\Length}{\upoperatorname{Length}}
\newcommand{\length}{\upoperatorname{length}}

\newcommand{\disc}{\upoperatorname{disc}}

\newcommand{\gr}{\upoperatorname{gr}}

\newcommand{\pr}{\upoperatorname{pr}}

\newcommand{\Frob}{\upoperatorname{Frob}}
\newcommand{\Gal}{\upoperatorname{Gal}}

\newcommand{\acts}{\curvearrowright}
\newcommand{\racts}{\curvearrowleft}

\renewcommand{\exp}{\upoperatorname{exp}}
\renewcommand{\log}{\upoperatorname{log}}
\renewcommand{\sin}{\upoperatorname{sin}}
\renewcommand{\cos}{\upoperatorname{cos}}
\renewcommand{\tan}{\upoperatorname{tan}}
\renewcommand{\arcsin}{\upoperatorname{arcsin}}
\renewcommand{\arccos}{\upoperatorname{arccos}}
\renewcommand{\arctan}{\upoperatorname{arctan}}

\newcommand{\Nm}{\upoperatorname{Nm}}



% for algebraic geometry
\newcommand{\Spec}{\upoperatorname{Spec}}
\newcommand{\Proj}{\upoperatorname{Proj}}
\newcommand{\mSpec}{\upoperatorname{mSpec}}
\newcommand{\Spf}{\upoperatorname{Spf}}

\newcommand{\calQuot}{\caloperatorname{Quot}}
\newcommand{\frakQuot}{\frakoperatorname{Quot}}
\newcommand{\Quot}{\upoperatorname{Quot}}
\newcommand{\Hilb}{\upoperatorname{Hilb}}
\newcommand{\calHilb}{\caloperatorname{Hilb}}
\newcommand{\frakHilb}{\frakoperatorname{Hilb}}
\newcommand{\Grass}{\upoperatorname{Grass}}
\newcommand{\calGrass}{\caloperatorname{Grass}}
\newcommand{\frakGrass}{\frakoperatorname{Grass}}
\newcommand{\frakMor}{\frakoperatorname{Mor}}


\newcommand{\calHom}{\caloperatorname{Hom}}
\newcommand{\calExt}{\caloperatorname{Ext}}
\newcommand{\calProj}{\caloperatorname{Proj}}
\newcommand{\calSpec}{\caloperatorname{Spec}}

\newcommand{\sing}{{\mathrm{sing}}}
\newcommand{\et}{{\text{\'et}}}
\newcommand{\zar}{{\mathrm{zar}}}
\newcommand{\nor}{{\mathrm{nor}}}

\newcommand{\CH}{\upoperatorname{CH}}

\newcommand{\ch}{\upoperatorname{ch}}
\newcommand{\td}{\upoperatorname{td}}


% for category theory and homological algebra
\newcommand{\Set}{\notationmathbold{Set}}
\newcommand{\Var}{\notationmathbold{Var}}
\newcommand{\Sch}{\notationmathbold{Sch}}
\newcommand{\sch}{\notationmathbold{sch}}
\newcommand{\Av}{\notationmathbold{Av}}
\newcommand{\Alg}{\notationmathbold{Alg}}
\newcommand{\Ring}{\notationmathbold{Ring}}
\newcommand{\Mod}{\notationmathbold{Mod}}
\newcommand{\module}{\notationmathbold{mod}}
\newcommand{\TopMod}{\notationmathbold{TopMod}}
\newcommand{\Banmod}{\notationmathbold{Banmod}}
\newcommand{\Vect}{\notationmathbold{Vect}}
\newcommand{\vect}{\notationmathbold{vect}}
\newcommand{\Grp}{\notationmathbold{Grp}}
\newcommand{\Grpd}{\notationmathbold{Grpd}}
\newcommand{\Ab}{\notationmathbold{Ab}}
\newcommand{\Top}{\notationmathbold{Top}}
\newcommand{\Open}{\notationmathbold{Open}}
\newcommand{\Sh}{\notationmathbold{Sh}}
\newcommand{\PSh}{\notationmathbold{PSh}}
\newcommand{\Coh}{\notationmathbold{Coh}}
\newcommand{\QCoh}{\notationmathbold{QCoh}}
\newcommand{\Cat}{\notationmathbold{Cat}}
\newcommand{\AlgSp}{\notationmathbold{AlgSp}}
\newcommand{\AlgGrp}{\notationmathbold{AlgGrp}}

\newcommand{\Obj}{\upoperatorname{Obj}}
\newcommand{\Fun}{\upoperatorname{Fun}}
\newcommand{\Mor}{\upoperatorname{Mor}}

\newcommand{\op}{\mathrm{op}}

\newcommand{\res}{\upoperatorname{res}}

\newcommand{\Ext}{\upoperatorname{Ext}}
\newcommand{\Tor}{\upoperatorname{Tor}}
\newcommand{\Hom}{\upoperatorname{Hom}}
\newcommand{\Eq}{\upoperatorname{Eq}}

\newcommand{\cohdim}{\upoperatorname{coh.dim}}
\newcommand{\hldim}{\upoperatorname{hl.dim}}
\newcommand{\projdim}{\upoperatorname{proj.dim}}
\newcommand{\injdim}{\upoperatorname{inj.dim}}



% for commutative algebra
\newcommand{\depth}{\upoperatorname{depth}}
\newcommand{\idealht}{\upoperatorname{ht}}
\newcommand{\codim}{\upoperatorname{codim}}
\newcommand{\Supp}{\upoperatorname{Supp}}
\newcommand{\trdeg}{\upoperatorname{trdeg}}
\newcommand{\Ass}{\upoperatorname{Ass}}
\newcommand{\Ann}{\upoperatorname{Ann}}
\newcommand{\Der}{\upoperatorname{Der}}

\newcommand{\Frac}{\upoperatorname{Frac}}

\newcommand{\Mult}{\upoperatorname{Mult}}

\newcommand{\rad}{\upoperatorname{rad}}
\newcommand{\nil}{\upoperatorname{nil}}
\newcommand{\characteristic}{\upoperatorname{char}}



% for birational geometry
\newcommand{\Pic}{\upoperatorname{Pic}}
\newcommand{\NS}{\upoperatorname{NS}}
\newcommand{\Exc}{\upoperatorname{Exc}}
\newcommand{\Sing}{\upoperatorname{Sing}}
\newcommand{\Bl}{\upoperatorname{Bl}}
\newcommand{\Reg}{\upoperatorname{Reg}}
\newcommand{\Psef}{\upoperatorname{Psef}}
\newcommand{\Nef}{\upoperatorname{Nef}}
\newcommand{\Amp}{\upoperatorname{Amp}}
\newcommand{\Bigcone}{\upoperatorname{Big}}
\newcommand{\Bs}{\upoperatorname{Bs}}
\newcommand{\Stab}{\upoperatorname{Stab}}
\newcommand{\Cone}{\upoperatorname{Cone}}
\newcommand{\Div}{\upoperatorname{Div}}
\newcommand{\WDiv}{\upoperatorname{WDiv}}
\newcommand{\CaDiv}{\upoperatorname{CaDiv}}
\newcommand{\divisor}{\upoperatorname{div}}
\newcommand{\Cl}{\upoperatorname{Cl}}
\newcommand{\CaCl}{\upoperatorname{CaCl}}
\newcommand{\WCl}{\upoperatorname{WCl}}
\newcommand{\ratmap}{\dashrightarrow}

\newcommand{\Bir}{\upoperatorname{Bir}}

\newcommand{\Bpf}{\upoperatorname{Bpf}}
\newcommand{\BPF}{\upoperatorname{BPF}}
\newcommand{\Mov}{\upoperatorname{Mov}}
\newcommand{\Nci}{\upoperatorname{Nci}}



% for Lie groups, Lie algebras and linear algebra
\newcommand{\GL}{\upoperatorname{GL}}
\newcommand{\PGL}{\upoperatorname{PGL}}
\newcommand{\SL}{\upoperatorname{SL}}
\newcommand{\PSL}{\upoperatorname{PSL}}
\newcommand{\SO}{\upoperatorname{SO}}
\newcommand{\Spin}{\upoperatorname{Spin}}
\newcommand{\Sp}{\upoperatorname{Sp}}
\newcommand{\SU}{\upoperatorname{SU}}

\newcommand{\Gr}{\upoperatorname{Gr}}

\newcommand{\gl}{{\mathfrak{gl}}}
\newcommand{\so}{{\mathfrak{so}}}
\newcommand{\su}{{\mathfrak{su}}}

\newcommand{\ad}{\upoperatorname{ad}}
\newcommand{\Ad}{\upoperatorname{Ad}}
\newcommand{\Lie}{\upoperatorname{Lie}}
\newcommand{\diag}{\upoperatorname{diag}}
\newcommand{\Span}{\upoperatorname{Span}}

\newcommand{\Image}{\upoperatorname{Im}}
\renewcommand{\Im}{\upoperatorname{Im}}
\newcommand{\Ker}{\upoperatorname{Ker}}
\renewcommand{\ker}{\upoperatorname{ker}}
\newcommand{\coker}{\upoperatorname{coker}}
\newcommand{\ord}{\upoperatorname{ord}}
\newcommand{\rank}{\upoperatorname{rank}}
\renewcommand{\deg}{\upoperatorname{deg}}
\renewcommand{\det}{\upoperatorname{det}}
\renewcommand{\dim}{\upoperatorname{dim}}
\newcommand{\tr}{\upoperatorname{tr}}
\newcommand{\Tr}{\upoperatorname{Tr}}

\newcommand{\Aut}{\upoperatorname{Aut}}
\newcommand{\End}{\upoperatorname{End}}
\newcommand{\Sym}{\upoperatorname{Sym}}

\newcommand{\MixDet}{\upoperatorname{MixDet}}




% analysis on cones
\newcommand{\Sep}{\upoperatorname{Sep}}




% for differential geometry and topology
\newcommand{\sm}{\mathrm{sm}}
\newcommand{\dR}{\mathrm{dR}}
\newcommand{\diff}{\mathrm{diff}}
\newcommand{\hol}{\mathrm{hol}}
\newcommand{\cnt}{\mathrm{cnt}}

\newcommand{\Hol}{\upoperatorname{Hol}}
\newcommand{\Mer}{\upoperatorname{Mer}}




% for Berkovich spaces
\newcommand{\Qnil}{\upoperatorname{Qnil}}
\newcommand{\Red}{\upoperatorname{Red}}
\newcommand{\cten}{\widehat{\otimes}}
\newcommand{\an}{\mathrm{an}}
\newcommand{\Int}{\upoperatorname{Int}}




% Arakelov geometry
\newcommand{\adeg}{\upoperatorname{\widehat{deg}}}
\newcommand{\aPic}{\upoperatorname{\widehat{Pic}}}
\newcommand{\aDiv}{\upoperatorname{\widehat{Div}}}
\newcommand{\adiv}{\upoperatorname{\widehat{divisor}}}
\newcommand{\aCaDiv}{\upoperatorname{\widehat{CaDiv}}}
\newcommand{\aWDiv}{\upoperatorname{\widehat{WDiv}}}
\newcommand{\aCl}{\upoperatorname{\widehat{Cl}}}
\newcommand{\aCaCl}{\upoperatorname{\widehat{CaCl}}}




% Class field theory
\newcommand{\unr}{{\text{unr}}}
\newcommand{\al}{{\text{al}}}
\newcommand{\ab}{{\text{ab}}}
\newcommand{\rec}{{\upoperatorname{rec}}}







\newcommand{\contradiction}{
	\raisebox{-0.6ex}{\makebox[2.4ex][c]{\begin{tikzpicture}
				\draw (0,0) circle (1.2ex);
				\draw (0.7 ex, 0.7 ex) -- (-0.4 ex, 0.7 ex);
				\draw (-0.4 ex, 0.7 ex) -- (0.7 ex, -0.4 ex);
				\draw (0.7 ex, -0.4 ex) -- (0.7 ex, 0.7 ex);
				\draw (0.7 ex, 0.7 ex) -- (-0.848 ex, -0.848 ex);
			\end{tikzpicture}}}
	\ \
}

% legendre符号
\makeatletter
\def\legendre@dash#1#2{\hb@xt@#1{%
		\kern-#2\p@
		\cleaders\hbox{\kern.5\p@
			\vrule\@height.2\p@\@depth.2\p@\@width\p@
			\kern.5\p@}\hfil
		\kern-#2\p@
	}}
\def\@legendre#1#2#3#4#5{\mathopen{}\left(
	\sbox\z@{$\genfrac{}{}{0pt}{#1}{#3#4}{#3#5}$}%
	\dimen@=\wd\z@
	\kern-\p@\vcenter{\box0}\kern-\dimen@\vcenter{\legendre@dash\dimen@{#2}}\kern-\p@
	\right)\mathclose{}}
\newcommand\legendre[2]{\mathchoice
	{\@legendre{0}{1}{}{#1}{#2}}
	{\@legendre{1}{.5}{\vphantom{1}}{#1}{#2}}
	{\@legendre{2}{0}{\vphantom{1}}{#1}{#2}}
	{\@legendre{3}{0}{\vphantom{1}}{#1}{#2}}
}
\def\dlegendre{\@legendre{0}{1}{}}
\def\tlegendre{\@legendre{1}{0.5}{\vphantom{1}}}
\makeatother





% 单一字母in mathbb mathcal mathsf mathrm mathfrak mathscr symbf symrm symup
\newcommand{\bbA}{{\mathbb{A}}}
\newcommand{\bba}{{\notationmathbb{a}}}
\newcommand{\bbB}{{\mathbb{B}}}
\newcommand{\bbb}{{\notationmathbb{b}}}
\newcommand{\bbC}{{\mathbb{C}}}
\newcommand{\bbc}{{\notationmathbb{c}}}
\newcommand{\bbD}{{\mathbb{D}}}
\newcommand{\bbd}{{\notationmathbb{d}}}
\newcommand{\bbE}{{\mathbb{E}}}
\newcommand{\bbe}{{\notationmathbb{e}}}
\newcommand{\bbF}{{\mathbb{F}}}
\newcommand{\bbf}{{\notationmathbb{f}}}
\newcommand{\bbG}{{\mathbb{G}}}
\newcommand{\bbg}{{\notationmathbb{g}}}
\newcommand{\bbH}{{\mathbb{H}}}
\newcommand{\bbh}{{\notationmathbb{h}}}
\newcommand{\bbI}{{\mathbb{I}}}
\newcommand{\bbi}{{\notationmathbb{i}}}
\newcommand{\bbJ}{{\mathbb{J}}}
\newcommand{\bbj}{{\notationmathbb{j}}}
\newcommand{\bbK}{{\mathbb{K}}}
\newcommand{\bbk}{{\notationmathbb{k}}}
\newcommand{\bbL}{{\mathbb{L}}}
\newcommand{\bbl}{{\notationmathbb{l}}}
\newcommand{\bbM}{{\mathbb{M}}}
\newcommand{\bbm}{{\notationmathbb{m}}}
\newcommand{\bbN}{{\mathbb{N}}}
\newcommand{\bbn}{{\notationmathbb{n}}}
\newcommand{\bbO}{{\mathbb{O}}}
\newcommand{\bbo}{{\notationmathbb{o}}}
\newcommand{\bbP}{{\mathbb{P}}}
\newcommand{\bbp}{{\notationmathbb{p}}}
\newcommand{\bbQ}{{\mathbb{Q}}}
\newcommand{\bbq}{{\notationmathbb{q}}}
\newcommand{\bbR}{{\mathbb{R}}}
\newcommand{\bbr}{{\notationmathbb{r}}}
\newcommand{\bbS}{{\mathbb{S}}}
\newcommand{\bbs}{{\notationmathbb{s}}}
\newcommand{\bbT}{{\mathbb{T}}}
\newcommand{\bbt}{{\notationmathbb{t}}}
\newcommand{\bbU}{{\mathbb{U}}}
\newcommand{\bbu}{{\notationmathbb{u}}}
\newcommand{\bbV}{{\mathbb{V}}}
\newcommand{\bbv}{{\notationmathbb{v}}}
\newcommand{\bbW}{{\mathbb{W}}}
\newcommand{\bbw}{{\notationmathbb{w}}}
\newcommand{\bbX}{{\mathbb{X}}}
\newcommand{\bbx}{{\notationmathbb{x}}}
\newcommand{\bbY}{{\mathbb{Y}}}
\newcommand{\bby}{{\notationmathbb{y}}}
\newcommand{\bbZ}{{\mathbb{Z}}}
\newcommand{\bbz}{{\notationmathbb{z}}}
\newcommand{\bbone}{{\notationmathbb{1}}}

\newcommand{\calA}{{\mathcal{A}}}
\newcommand{\cala}{{\notationmathcal{a}}}
\newcommand{\calB}{{\mathcal{B}}}
\newcommand{\calb}{{\notationmathcal{b}}}
\newcommand{\calC}{{\mathcal{C}}}
\newcommand{\calc}{{\notationmathcal{c}}}
\newcommand{\calD}{{\mathcal{D}}}
\newcommand{\cald}{{\notationmathcal{d}}}
\newcommand{\calE}{{\mathcal{E}}}
\newcommand{\cale}{{\notationmathcal{e}}}
\newcommand{\calF}{{\mathcal{F}}}
\newcommand{\calf}{{\notationmathcal{f}}}
\newcommand{\calG}{{\mathcal{G}}}
\newcommand{\calg}{{\notationmathcal{g}}}
\newcommand{\calH}{{\mathcal{H}}}
\newcommand{\calh}{{\notationmathcal{h}}}
\newcommand{\calI}{{\mathcal{I}}}
\newcommand{\cali}{{\notationmathcal{i}}}
\newcommand{\calJ}{{\mathcal{J}}}
\newcommand{\calj}{{\notationmathcal{j}}}
\newcommand{\calK}{{\mathcal{K}}}
\newcommand{\calk}{{\notationmathcal{k}}}
\newcommand{\calL}{{\mathcal{L}}}
\newcommand{\call}{{\notationmathcal{l}}}
\newcommand{\calM}{{\mathcal{M}}}
\newcommand{\calm}{{\notationmathcal{m}}}
\newcommand{\calN}{{\mathcal{N}}}
\newcommand{\caln}{{\notationmathcal{n}}}
\newcommand{\calO}{{\mathcal{O}}}
\newcommand{\calo}{{\notationmathcal{o}}}
\newcommand{\calP}{{\mathcal{P}}}
\newcommand{\calp}{{\notationmathcal{p}}}
\newcommand{\calQ}{{\mathcal{Q}}}
\newcommand{\calq}{{\notationmathcal{q}}}
\newcommand{\calR}{{\mathcal{R}}}
\newcommand{\calr}{{\notationmathcal{r}}}
\newcommand{\calS}{{\mathcal{S}}}
\newcommand{\cals}{{\notationmathcal{s}}}
\newcommand{\calT}{{\mathcal{T}}}
\newcommand{\calt}{{\notationmathcal{t}}}
\newcommand{\calU}{{\mathcal{U}}}
\newcommand{\calu}{{\notationmathcal{u}}}
\newcommand{\calV}{{\mathcal{V}}}
\newcommand{\calv}{{\notationmathcal{v}}}
\newcommand{\calW}{{\mathcal{W}}}
\newcommand{\calw}{{\notationmathcal{w}}}
\newcommand{\calX}{{\mathcal{X}}}
\newcommand{\calx}{{\notationmathcal{x}}}
\newcommand{\calY}{{\mathcal{Y}}}
\newcommand{\caly}{{\notationmathcal{y}}}
\newcommand{\calZ}{{\mathcal{Z}}}
\newcommand{\calz}{{\notationmathcal{z}}}

\newcommand{\scrA}{{\mathscr{A}}}
\newcommand{\scra}{{\notationmathscr{a}}}
\newcommand{\scrB}{{\mathscr{B}}}
\newcommand{\scrb}{{\notationmathscr{b}}}
\newcommand{\scrC}{{\mathscr{C}}}
\newcommand{\scrc}{{\notationmathscr{c}}}
\newcommand{\scrD}{{\mathscr{D}}}
\newcommand{\scrd}{{\notationmathscr{d}}}
\newcommand{\scrE}{{\mathscr{E}}}
\newcommand{\scre}{{\notationmathscr{e}}}
\newcommand{\scrF}{{\mathscr{F}}}
\newcommand{\scrf}{{\notationmathscr{f}}}
\newcommand{\scrG}{{\mathscr{G}}}
\newcommand{\scrg}{{\notationmathscr{g}}}
\newcommand{\scrH}{{\mathscr{H}}}
\newcommand{\scrh}{{\notationmathscr{h}}}
\newcommand{\scrI}{{\mathscr{I}}}
\newcommand{\scri}{{\notationmathscr{i}}}
\newcommand{\scrJ}{{\mathscr{J}}}
\newcommand{\scrj}{{\notationmathscr{j}}}
\newcommand{\scrK}{{\mathscr{K}}}
\newcommand{\scrk}{{\notationmathscr{k}}}
\newcommand{\scrL}{{\mathscr{L}}}
\newcommand{\scrl}{{\notationmathscr{l}}}
\newcommand{\scrM}{{\mathscr{M}}}
\newcommand{\scrm}{{\notationmathscr{m}}}
\newcommand{\scrN}{{\mathscr{N}}}
\newcommand{\scrn}{{\notationmathscr{n}}}
\newcommand{\scrO}{{\mathscr{O}}}
\newcommand{\scro}{{\notationmathscr{o}}}
\newcommand{\scrP}{{\mathscr{P}}}
\newcommand{\scrp}{{\notationmathscr{p}}}
\newcommand{\scrQ}{{\mathscr{Q}}}
\newcommand{\scrq}{{\notationmathscr{q}}}
\newcommand{\scrR}{{\mathscr{R}}}
\newcommand{\scrr}{{\notationmathscr{r}}}
\newcommand{\scrS}{{\mathscr{S}}}
\newcommand{\scrs}{{\notationmathscr{s}}}
\newcommand{\scrT}{{\mathscr{T}}}
\newcommand{\scrt}{{\notationmathscr{t}}}
\newcommand{\scrU}{{\mathscr{U}}}
\newcommand{\scru}{{\notationmathscr{u}}}
\newcommand{\scrV}{{\mathscr{V}}}
\newcommand{\scrv}{{\notationmathscr{v}}}
\newcommand{\scrW}{{\mathscr{W}}}
\newcommand{\scrw}{{\notationmathscr{w}}}
\newcommand{\scrX}{{\mathscr{X}}}
\newcommand{\scrx}{{\notationmathscr{x}}}
\newcommand{\scrY}{{\mathscr{Y}}}
\newcommand{\scry}{{\notationmathscr{y}}}
\newcommand{\scrZ}{{\mathscr{Z}}}
\newcommand{\scrz}{{\notationmathscr{z}}}

\newcommand{\frakA}{{\mathfrak{A}}}
\newcommand{\fraka}{{\mathfrak{a}}}
\newcommand{\frakB}{{\mathfrak{B}}}
\newcommand{\frakb}{{\mathfrak{b}}}
\newcommand{\frakC}{{\mathfrak{C}}}
\newcommand{\frakc}{{\mathfrak{c}}}
\newcommand{\frakD}{{\mathfrak{D}}}
\newcommand{\frakd}{{\mathfrak{d}}}
\newcommand{\frakE}{{\mathfrak{E}}}
\newcommand{\frake}{{\mathfrak{e}}}
\newcommand{\frakF}{{\mathfrak{F}}}
\newcommand{\frakf}{{\mathfrak{f}}}
\newcommand{\frakG}{{\mathfrak{G}}}
\newcommand{\frakg}{{\mathfrak{g}}}
\newcommand{\frakH}{{\mathfrak{H}}}
\newcommand{\frakh}{{\mathfrak{h}}}
\newcommand{\frakI}{{\mathfrak{I}}}
\newcommand{\fraki}{{\mathfrak{i}}}
\newcommand{\frakJ}{{\mathfrak{J}}}
\newcommand{\frakj}{{\mathfrak{j}}}
\newcommand{\frakK}{{\mathfrak{K}}}
\newcommand{\frakk}{{\mathfrak{k}}}
\newcommand{\frakL}{{\mathfrak{L}}}
\newcommand{\frakl}{{\mathfrak{l}}}
\newcommand{\frakM}{{\mathfrak{M}}}
\newcommand{\frakm}{{\mathfrak{m}}}
\newcommand{\frakN}{{\mathfrak{N}}}
\newcommand{\frakn}{{\mathfrak{n}}}
\newcommand{\frakO}{{\mathfrak{O}}}
\newcommand{\frako}{{\mathfrak{o}}}
\newcommand{\frakP}{{\mathfrak{P}}}
\newcommand{\frakp}{{\mathfrak{p}}}
\newcommand{\frakQ}{{\mathfrak{Q}}}
\newcommand{\frakq}{{\mathfrak{q}}}
\newcommand{\frakR}{{\mathfrak{R}}}
\newcommand{\frakr}{{\mathfrak{r}}}
\newcommand{\frakS}{{\mathfrak{S}}}
\newcommand{\fraks}{{\mathfrak{s}}}
\newcommand{\frakT}{{\mathfrak{T}}}
\newcommand{\frakt}{{\mathfrak{t}}}
\newcommand{\frakU}{{\mathfrak{U}}}
\newcommand{\fraku}{{\mathfrak{u}}}
\newcommand{\frakV}{{\mathfrak{V}}}
\newcommand{\frakv}{{\mathfrak{v}}}
\newcommand{\frakW}{{\mathfrak{W}}}
\newcommand{\frakw}{{\mathfrak{w}}}
\newcommand{\frakX}{{\mathfrak{X}}}
\newcommand{\frakx}{{\mathfrak{x}}}
\newcommand{\frakY}{{\mathfrak{Y}}}
\newcommand{\fraky}{{\mathfrak{y}}}
\newcommand{\frakZ}{{\mathfrak{Z}}}
\newcommand{\frakz}{{\mathfrak{z}}}

\newcommand{\rmA}{{\notationmathroman{A}}}
\newcommand{\rma}{{\notationmathroman{a}}}
\newcommand{\rmB}{{\notationmathroman{B}}}
\newcommand{\rmb}{{\notationmathroman{b}}}
\newcommand{\rmC}{{\notationmathroman{C}}}
\newcommand{\rmc}{{\notationmathroman{c}}}
\newcommand{\rmD}{{\notationmathroman{D}}}
\newcommand{\rmd}{{\notationmathroman{d}}}
\newcommand{\rmE}{{\notationmathroman{E}}}
\newcommand{\rme}{{\notationmathroman{e}}}
\newcommand{\rmF}{{\notationmathroman{F}}}
\newcommand{\rmf}{{\notationmathroman{f}}}
\newcommand{\rmG}{{\notationmathroman{G}}}
\newcommand{\rmg}{{\notationmathroman{g}}}
\newcommand{\rmH}{{\notationmathroman{H}}}
\newcommand{\rmh}{{\notationmathroman{h}}}
\newcommand{\rmI}{{\notationmathroman{I}}}
\newcommand{\rmi}{{\notationmathroman{i}}}
\newcommand{\rmJ}{{\notationmathroman{J}}}
\newcommand{\rmj}{{\notationmathroman{j}}}
\newcommand{\rmK}{{\notationmathroman{K}}}
\newcommand{\rmk}{{\notationmathroman{k}}}
\newcommand{\rmL}{{\notationmathroman{L}}}
\newcommand{\rml}{{\notationmathroman{l}}}
\newcommand{\rmM}{{\notationmathroman{M}}}
\newcommand{\rmm}{{\notationmathroman{m}}}
\newcommand{\rmN}{{\notationmathroman{N}}}
\newcommand{\rmn}{{\notationmathroman{n}}}
\newcommand{\rmO}{{\notationmathroman{O}}}
\newcommand{\rmo}{{\notationmathroman{o}}}
\newcommand{\rmP}{{\notationmathroman{P}}}
\newcommand{\rmp}{{\notationmathroman{p}}}
\newcommand{\rmQ}{{\notationmathroman{Q}}}
\newcommand{\rmq}{{\notationmathroman{q}}}
\newcommand{\rmR}{{\notationmathroman{R}}}
\newcommand{\rmr}{{\notationmathroman{r}}}
\newcommand{\rmS}{{\notationmathroman{S}}}
\newcommand{\rms}{{\notationmathroman{s}}}
\newcommand{\rmT}{{\notationmathroman{T}}}
\newcommand{\rmt}{{\notationmathroman{t}}}
\newcommand{\rmU}{{\notationmathroman{U}}}
\newcommand{\rmu}{{\notationmathroman{u}}}
\newcommand{\rmV}{{\notationmathroman{V}}}
\newcommand{\rmv}{{\notationmathroman{v}}}
\newcommand{\rmW}{{\notationmathroman{W}}}
\newcommand{\rmw}{{\notationmathroman{w}}}
\newcommand{\rmX}{{\notationmathroman{X}}}
\newcommand{\rmx}{{\notationmathroman{x}}}
\newcommand{\rmY}{{\notationmathroman{Y}}}
\newcommand{\rmy}{{\notationmathroman{y}}}
\newcommand{\rmZ}{{\notationmathroman{Z}}}
\newcommand{\rmz}{{\notationmathroman{z}}}

\newcommand{\upA}{{\notationmathupright{A}}}
\newcommand{\upa}{{\notationmathupright{a}}}
\newcommand{\upB}{{\notationmathupright{B}}}
\newcommand{\upb}{{\notationmathupright{b}}}
\newcommand{\upC}{{\notationmathupright{C}}}
\newcommand{\upc}{{\notationmathupright{c}}}
\newcommand{\upD}{{\notationmathupright{D}}}
\newcommand{\upd}{{\notationmathupright{d}}}
\newcommand{\upE}{{\notationmathupright{E}}}
\newcommand{\upe}{{\notationmathupright{e}}}
\newcommand{\upF}{{\notationmathupright{F}}}
\newcommand{\upf}{{\notationmathupright{f}}}
\newcommand{\upG}{{\notationmathupright{G}}}
\newcommand{\upg}{{\notationmathupright{g}}}
\newcommand{\upH}{{\notationmathupright{H}}}
\newcommand{\uph}{{\notationmathupright{h}}}
\newcommand{\upI}{{\notationmathupright{I}}}
\newcommand{\upi}{{\notationmathupright{i}}}
\newcommand{\upJ}{{\notationmathupright{J}}}
\newcommand{\upj}{{\notationmathupright{j}}}
\newcommand{\upK}{{\notationmathupright{K}}}
\newcommand{\upk}{{\notationmathupright{k}}}
\newcommand{\upL}{{\notationmathupright{L}}}
\newcommand{\upl}{{\notationmathupright{l}}}
\newcommand{\upM}{{\notationmathupright{M}}}
\newcommand{\upm}{{\notationmathupright{m}}}
\newcommand{\upN}{{\notationmathupright{N}}}
\newcommand{\upn}{{\notationmathupright{n}}}
\newcommand{\upO}{{\notationmathupright{O}}}
\newcommand{\upo}{{\notationmathupright{o}}}
\newcommand{\upP}{{\notationmathupright{P}}}
\newcommand{\upp}{{\notationmathupright{p}}}
\newcommand{\upQ}{{\notationmathupright{Q}}}
\newcommand{\upq}{{\notationmathupright{q}}}
\newcommand{\upR}{{\notationmathupright{R}}}
\newcommand{\upr}{{\notationmathupright{r}}}
\newcommand{\upS}{{\notationmathupright{S}}}
\newcommand{\ups}{{\notationmathupright{s}}}
\newcommand{\upT}{{\notationmathupright{T}}}
\newcommand{\upt}{{\notationmathupright{t}}}
\newcommand{\upU}{{\notationmathupright{U}}}
\newcommand{\upu}{{\notationmathupright{u}}}
\newcommand{\upV}{{\notationmathupright{V}}}
\newcommand{\upv}{{\notationmathupright{v}}}
\newcommand{\upW}{{\notationmathupright{W}}}
\newcommand{\upw}{{\notationmathupright{w}}}
\newcommand{\upX}{{\notationmathupright{X}}}
\newcommand{\upx}{{\notationmathupright{x}}}
\newcommand{\upY}{{\notationmathupright{Y}}}
\newcommand{\upy}{{\notationmathupright{y}}}
\newcommand{\upZ}{{\notationmathupright{Z}}}
\newcommand{\upz}{{\notationmathupright{z}}}

\newcommand{\bfA}{{\notationmathbold{A}}}
\newcommand{\bfa}{{\notationmathbold{a}}}
\newcommand{\bfB}{{\notationmathbold{B}}}
\newcommand{\bfb}{{\notationmathbold{b}}}
\newcommand{\bfC}{{\notationmathbold{C}}}
\newcommand{\bfc}{{\notationmathbold{c}}}
\newcommand{\bfD}{{\notationmathbold{D}}}
\newcommand{\bfd}{{\notationmathbold{d}}}
\newcommand{\bfE}{{\notationmathbold{E}}}
\newcommand{\bfe}{{\notationmathbold{e}}}
\newcommand{\bfF}{{\notationmathbold{F}}}
\newcommand{\bff}{{\notationmathbold{f}}}
\newcommand{\bfG}{{\notationmathbold{G}}}
\newcommand{\bfg}{{\notationmathbold{g}}}
\newcommand{\bfH}{{\notationmathbold{H}}}
\newcommand{\bfh}{{\notationmathbold{h}}}
\newcommand{\bfI}{{\notationmathbold{I}}}
\newcommand{\bfi}{{\notationmathbold{i}}}
\newcommand{\bfJ}{{\notationmathbold{J}}}
\newcommand{\bfj}{{\notationmathbold{j}}}
\newcommand{\bfK}{{\notationmathbold{K}}}
\newcommand{\bfk}{{\notationmathbold{k}}}
\newcommand{\bfL}{{\notationmathbold{L}}}
\newcommand{\bfl}{{\notationmathbold{l}}}
\newcommand{\bfM}{{\notationmathbold{M}}}
\newcommand{\bfm}{{\notationmathbold{m}}}
\newcommand{\bfN}{{\notationmathbold{N}}}
\newcommand{\bfn}{{\notationmathbold{n}}}
\newcommand{\bfO}{{\notationmathbold{O}}}
\newcommand{\bfo}{{\notationmathbold{o}}}
\newcommand{\bfP}{{\notationmathbold{P}}}
\newcommand{\bfp}{{\notationmathbold{p}}}
\newcommand{\bfQ}{{\notationmathbold{Q}}}
\newcommand{\bfq}{{\notationmathbold{q}}}
\newcommand{\bfR}{{\notationmathbold{R}}}
\newcommand{\bfr}{{\notationmathbold{r}}}
\newcommand{\bfS}{{\notationmathbold{S}}}
\newcommand{\bfs}{{\notationmathbold{s}}}
\newcommand{\bfT}{{\notationmathbold{T}}}
\newcommand{\bft}{{\notationmathbold{t}}}
\newcommand{\bfU}{{\notationmathbold{U}}}
\newcommand{\bfu}{{\notationmathbold{u}}}
\newcommand{\bfV}{{\notationmathbold{V}}}
\newcommand{\bfv}{{\notationmathbold{v}}}
\newcommand{\bfW}{{\notationmathbold{W}}}
\newcommand{\bfw}{{\notationmathbold{w}}}
\newcommand{\bfX}{{\notationmathbold{X}}}
\newcommand{\bfx}{{\notationmathbold{x}}}
\newcommand{\bfY}{{\notationmathbold{Y}}}
\newcommand{\bfy}{{\notationmathbold{y}}}
\newcommand{\bfZ}{{\notationmathbold{Z}}}
\newcommand{\bfz}{{\notationmathbold{z}}}

\newcommand{\sfA}{{\mathsf{A}}}
\newcommand{\sfa}{{\mathsf{a}}}
\newcommand{\sfB}{{\mathsf{B}}}
\newcommand{\sfb}{{\mathsf{b}}}
\newcommand{\sfC}{{\mathsf{C}}}
\newcommand{\sfc}{{\mathsf{c}}}
\newcommand{\sfD}{{\mathsf{D}}}
\newcommand{\sfd}{{\mathsf{d}}}
\newcommand{\sfE}{{\mathsf{E}}}
\newcommand{\sfe}{{\mathsf{e}}}
\newcommand{\sfF}{{\mathsf{F}}}
\newcommand{\sff}{{\mathsf{f}}}
\newcommand{\sfG}{{\mathsf{G}}}
\newcommand{\sfg}{{\mathsf{g}}}
\newcommand{\sfH}{{\mathsf{H}}}
\newcommand{\sfh}{{\mathsf{h}}}
\newcommand{\sfI}{{\mathsf{I}}}
\newcommand{\sfi}{{\mathsf{i}}}
\newcommand{\sfJ}{{\mathsf{J}}}
\newcommand{\sfj}{{\mathsf{j}}}
\newcommand{\sfK}{{\mathsf{K}}}
\newcommand{\sfk}{{\mathsf{k}}}
\newcommand{\sfL}{{\mathsf{L}}}
\newcommand{\sfl}{{\mathsf{l}}}
\newcommand{\sfM}{{\mathsf{M}}}
\newcommand{\sfm}{{\mathsf{m}}}
\newcommand{\sfN}{{\mathsf{N}}}
\newcommand{\sfn}{{\mathsf{n}}}
\newcommand{\sfO}{{\mathsf{O}}}
\newcommand{\sfo}{{\mathsf{o}}}
\newcommand{\sfP}{{\mathsf{P}}}
\newcommand{\sfp}{{\mathsf{p}}}
\newcommand{\sfQ}{{\mathsf{Q}}}
\newcommand{\sfq}{{\mathsf{q}}}
\newcommand{\sfR}{{\mathsf{R}}}
\newcommand{\sfr}{{\mathsf{r}}}
\newcommand{\sfS}{{\mathsf{S}}}
\newcommand{\sfs}{{\mathsf{s}}}
\newcommand{\sfT}{{\mathsf{T}}}
\newcommand{\sft}{{\mathsf{t}}}
\newcommand{\sfU}{{\mathsf{U}}}
\newcommand{\sfu}{{\mathsf{u}}}
\newcommand{\sfV}{{\mathsf{V}}}
\newcommand{\sfv}{{\mathsf{v}}}
\newcommand{\sfW}{{\mathsf{W}}}
\newcommand{\sfw}{{\mathsf{w}}}
\newcommand{\sfX}{{\mathsf{X}}}
\newcommand{\sfx}{{\mathsf{x}}}
\newcommand{\sfY}{{\mathsf{Y}}}
\newcommand{\sfy}{{\mathsf{y}}}
\newcommand{\sfZ}{{\mathsf{Z}}}
\newcommand{\sfz}{{\mathsf{z}}}

\usetikzlibrary{arrows.meta,calc}

\def\acts{\ \rotatebox[origin=c]{-90}{$\circlearrowright$}\ }
\def\racts{\ \rotatebox[origin=c]{90}{$\circlearrowleft$}\ }
% \def\dacts{\ \rotatebox[origin=c]{0}{$\circlearrowright$}\ }
% \def\uacts{\ \rotatebox[origin=c]{180}{$\circlearrowright$}\ }

\title{Algebraic Dynamics and Dynamical Iitaka Theory}
% \subtitle{}
\date{Undergraduate Forum, ICCM 2025, January 7, 2026}
\author{Tianle Yang \texorpdfstring{\\}{,} based on the joint work with Sheng Meng and Long Wang}


% Institute logo on the title page
\titlegraphic{\hfill\includegraphics[height=1.2cm]{\ECNUMathLogo}}


\begin{document}

% First page with different background
{
    \setbeamertemplate{background}{
        \includegraphics[width=\paperwidth,height=\paperheight]{\ECNUFirstPageBG}
    }
    \maketitle
}


% ----------------------------------------------------------


\begin{frame}{Kawaguchi-Silverman Conjecture}

    \vspace{0.3em}
        
    \(/\overline{\bbQ}\). \(X\) smooth projective variety; \(f:X \to X\) surjective endomorphism.

    \pause
    \vspace{0.2em}

    \begin{conjecture}[{Kawaguchi-Silverman Conjecture $=$ KSC}]\label{conj: ksc} 
        If the orbit \(O_f(x):=\{f^n(x)\,|\,n\ge 0\}\) is Zariski dense in \(X\), then
        \[ \alpha_f(x) = \delta_f.  \]
    \end{conjecture} 

    \pause
    \vspace{0.2em}

    Let \(H\) be an ample divisor and \(h\geq 1\) a height function associated to \(H\).

    % \vspace{0.5em}
    \begin{center}
    \begin{tabular}{cll}
        \visible<4->{
            \textbullet & \textcolor{magenta}{arithmetic degree}: & \\
                & \(\displaystyle \textcolor{magenta}{\alpha_f(x)} := \lim_{n\to\infty} h(f^n(x))^{\frac{1}{n}}\)
                & \textcolor{magenta}{arithmetic, local} invariant at \(x\); \\
        }
        \visible<5->{
            \textbullet & \textcolor{cyan}{dynamical degree}: & \\
                & \(\displaystyle \textcolor{cyan}{\delta_f} := \lim_{n\to\infty} ((f^n)^*H \cdot H^{\dim X - 1})^{\frac{1}{n}}\) 
                & \textcolor{cyan}{geometric, global} invariant of \(f\). \\
        }
    \end{tabular}
    \end{center}

    \vspace{-0.7cm}
    \visible<6->{
        \begin{center}
            \textbf{Slogan:} {\Large\textcolor{cyan}{GEOMETRY}} controls {\Large\textcolor{magenta}{ARITHMETIC}}.
        \end{center}
    }
    
\end{frame}


\begin{frame}{An example}
    \vskip-1em
    \begin{tikzpicture}
        % axes from -4 to 4 with grid and ticks
        \draw[help lines, gray!20] (-3.5,-2.5) grid (4.5,4.5);
        \draw[->, thick] (-3.5,0) -- (4.5,0) node[right] {$x$};
        \draw[->, thick] (0,-2.5) -- (0,4.5) node[above] {$y$};
        \foreach \pos/\labels in {-3/{-1000}, -2/{-100}, -1/{-10}, 1/{10}, 2/{100}, 3/{1000}, 4/{10000}}{
            \draw (\pos,0.06) -- (\pos,-0.06) node[below,text=gray!50,font=\scriptsize] {\labels};
        }
        \foreach \pos/\labels in {-2/{0.01}, -1/{0.1}, 1/{10}, 2/{100}, 3/{1000}, 4/{10000}}{
            \draw (0.06,\pos) -- (-0.06,\pos) node[left,text=gray!50,font=\scriptsize] {\labels};
        }

        % draw the orbit points and arrows between them
        
        % define named coordinates for the orbit points
        \coordinate (p0) at (0.05,0.301);
        \coordinate (p1) at (-0.477,0.301);
        \coordinate (p2) at (0.699,-0.778);
        \coordinate (p3) at (-1.041,-1.477);
        \coordinate (p4) at (-2.892,2.342);
        \coordinate (p5) at (5.747,5.234);

        % draw the filled points
        \foreach \p in {p0,p1,p2,p3,p4,p5}{
            \fill[red!80] (\p) circle (0.08);
        }

        % draw labels for the points
        \node[above right,font=\scriptsize,text=red!80] at (p0) {\([1:2:1]\)};
        \node[above left,font=\scriptsize,text=red!80] at (p1) {\([-3:2:1]\)};
        \node[below right,font=\scriptsize,text=red!80] at (p2) {\([5:-6:1]\)};
        \node[below left,font=\scriptsize,text=red!80] at (p3) {\([-11:-30:1]\)};
        \node[right,font=\scriptsize,text=red!80] at (p4) {\([-779:220:1]\)};
        \node[above right,font=\scriptsize,text=red!80] at (p5) {\([558441:171380:1]\)};

        % draw arrows between named points
        \draw[|->,thick,orange,shorten <=0.09cm,shorten >=0.09cm] (p0) -- (p1);
        \draw[|->,thick,orange,shorten <=0.09cm,shorten >=0.09cm] (p1) -- (p2);
        \draw[|->,thick,orange,shorten <=0.09cm,shorten >=0.09cm] (p2) -- (p3);
        \draw[|->,thick,orange,shorten <=0.09cm,shorten >=0.09cm] (p3) -- (p4);
        \draw[|->,thick,orange,shorten <=0.09cm,shorten >=0.09cm] (p4) -- (p5);

    \end{tikzpicture}

    \begin{tikzpicture}[remember picture,overlay]
    \node[anchor=south east,xshift=-0.6cm,yshift=0.8cm] at (current page.south east) {%
        \begin{minipage}[c][8cm]{0.44\textwidth}
            {\fontsize{9pt}{11pt}\selectfont
                \setbeamertemplate{itemize item}{\Large\textbullet}
                \begin{itemize}[<+- | alert@+>]
                    \setlength{\itemsep}{8pt}
                    \setlength{\parskip}{2pt}
                    \item 
                    \(X=\mathbb{P}^2\), 
                    
                    \(f:[x:y:z]\mapsto[x^2-y^2:xy:z^2]\),
                    
                    \(x = [1:2:1]\).

                    \item \(f^*H \sim 2H\) \(\Rightarrow\) \(\delta_f = 2\).
                    
                    \item 
                    \begin{tabular}{r l l}
                    \(n\) & \(h(f^n(x))\)   &   \\ \hline
                    \(0\) & \(\log 2\)      & \approx \(0.7\) \\
                    \(1\) & \(\log 3\)      & \approx \(1.1\) \\
                    \(2\) & \(\log 6\)      & \approx \(1.8\) \\
                    \(3\) & \(\log 30\)     & \approx \(3.4\) \\
                    \(4\) & \(\log 779\)    & \approx \(6.7\) \\
                    \(5\) & \(\log 558441\) & \approx \(13.2\)
                    \end{tabular}

                    \item It is expected that \(\alpha_f(x) = 2\).

                \end{itemize}
            }
        \end{minipage}
    };
    \end{tikzpicture}

    % Let \(X=\mathbb{P}^2\) and \(f:[x:y:z]\mapsto[x^2-y^2:xy:z^2]\).

    % \(x=[1:2:1]\) has a Zariski dense orbit and \(\delta_f=2\).
    % \(f(x)=[-3:2:1]\), 
    % \(f^2(x)=[5:-6:1]\), 
    % \(f^3(x)=[-11:-30:1]\), 
    % \(f^4(x)=[-779:220:1]\), 
    % \(f^5(x)=[558441:171380:1]\), 

    % log10 coordinates of the orbit points:
    % \((0,0.301),(-0.477,0.301),(0.699,-0.778),(-1.041,-1.477),(-2.892,2.342),(5.747,5.234),\ldots\)

    % Heights:
    % \(h(x)=\log\max\{|1|,|2|,|1|\}=\log2\approx0.7\),
    % \(h(f(x))=\log3\approx1.1\),
    % \(h(f^2(x))=\log6\approx1.8\),
    % \(h(f^3(x))=\log30\approx3.4\),
    % \(h(f^4(x))=\log779\approx6.7\),
    % \(h(f^5(x))=\log558441\approx13.2,\ldots\)
\end{frame}


\begin{frame}[c]{Three orbit conjectures}

    \begin{tikzpicture}[>=Stealth, every node/.style={font=\small}, scale=0.7, transform shape]
        % nodes
        \node[draw=blue!60, fill=blue!6, thick, circle, minimum size=2.5cm, align=center] (ZDO) at (-2.5,0) {Zariski Dense\\ Orbit conjecture\\ \textcolor{blue!70!black}{(ZDO)}};
        \node[draw=green!60, fill=green!6, thick, circle, minimum size=2.5cm, align=center] (DML) at (2.5,0) {Dynamical\\ Mordell-Lang\\ conjecture \textcolor{green!70!black}{(DML)}};
        \node[draw=red!70, fill=red!6, thick, circle, minimum size=2.5cm, align=center] (KSC) at (0,-3.7) {\textbf{Kawaguchi-}\\\textbf{Silverman}\\ \textbf{conjecture} \textcolor{red!70!black}{\textbf{(KSC)}}};

        % curved association arrows
        \draw[-, thick, orange] (ZDO) to[out=60,in=120] (DML);
        \draw[-, thick, orange] (DML) to[out=-60,in=-20] (KSC);
        \draw[-, thick, orange] (KSC) to[out=-160,in=-120] (ZDO);
    \end{tikzpicture}

    \begin{tikzpicture}[remember picture,overlay]
    \node[anchor=east,xshift=-0.6cm] at (current page.east) {%
        \begin{minipage}[t][0.9\textheight]{0.55\textwidth}
            \only<1-3>{
                \vspace{0.5cm}
                \textcolor{red!70!black}{
                    The KSC relates the growth rate of heights along orbits to dynamical degrees.
                }
            }

            \only<2-3>{
                \vspace{0.5cm}
                \textcolor{green!70!black}{
                    % The DML predicts if 
                    % \[ \# O_f(x) \cap V = \infty \quad V \text{ a subvariety}, \]
                    % then 
                    % \[ O_{f^r}(f^s(x)) \subseteq V \quad \text{for some } r,s. \]
                    The DML describes the intersection between orbits and subvarieties.
                }
            }

            \only<3>{
                \vspace{0.5cm}
                \textcolor{blue!70!black}{
                    The ZDO states that either 
                    \[ \exists x\  \text{with}\  \overline{O_f(x)} = X, \]
                    or \(f\) descends to identity after iteration.
                    % \[ \exists\  f^n \acts X \ratmap Y \racts \id_Y. \]
                }
            }

            \only<4->{Main known cases:}

            \only<5>{
                \vspace{0.3cm}
                {\fontsize{9pt}{11pt}\selectfont
                \textcolor{red!70!black}{Smooth projective surfaces}                        \hfill [Matsuzawa-Sano-Shibata];\\
                \textcolor{red!70!black}{Quasi-projective surfaces (assume DML)}            \hfill [Wang];\\
                \textcolor{red!70!black}{Smooth projective 3folds with \(\deg f>1\)}        \hfill [Meng-Zhang];\\
                \vspace{0.1em}\\
                \textcolor{red!70!black}{Abelian varieties}                                 \hfill [Kawaguchi-Silverman];\\
                \textcolor{red!70!black}{Hyperk\"ahler manifolds}                           \hfill [Lesieutre-Santriano];\\
                \textcolor{red!70!black}{Mori dream spaces}                                 \hfill [Matsuzawa];\\
                \vspace{0.1em}\\
                \textcolor{red!70!black}{Polarized endomorphisms}                           \hfill [Kawaguchi-Silverman];\\
                \textcolor{red!70!black}{Int-amplified endomorphisms}                       \hfill [Meng-Zhong];\\
                \vspace{0.1em}\\
                \textcolor{red!70!black}{\(\cdots\)}
                }
            }

            \only<6-7>{
                \vspace{0.3cm}
                {\fontsize{9pt}{11pt}\selectfont
                \textcolor{green!70!black}{\'Etale case}                            \hfill [Bell-Ghioca-Tucker];\\
                \textcolor{green!70!black}{Endomorphisms of \(\bbA^2\)}             \hfill [Xie];\\
                \textcolor{green!70!black}{ Projective surfaces in char \(p>0\)}    \hfill [Xie-Yang];\\
                \textcolor{green!70!black}{\(\cdots\)}
                }
            }

            \only<7>{
                \vspace{0.3cm}
                {\fontsize{9pt}{11pt}\selectfont
                \textcolor{blue!70!black}{Over \(\bbC\) (uncountable field)}                    \hfill [Amerik-Campana];\\
                \textcolor{blue!70!black}{Existence of infinity orbits}                         \hfill [Amerik];\\
                \textcolor{blue!70!black}{Projective surfaces}                                  \hfill [Xie,Jia-Xie-Zhang];\\
                \textcolor{blue!70!black}{Automorphisms of threefolds with positive entropy}    \hfill [Matsuzawa-Xie];\\
                \textcolor{blue!70!black}{\(\cdots\)}
                }
            }
        \end{minipage}
    };
    \end{tikzpicture}
    
\end{frame}


\begin{frame}{Main results}

    \only<1->{
        Settings:
        \begin{itemize}
            \item \(\pi:(X,f) \to (Y,g)\) smooth Fano fibration with \(X,Y\) smooth projective and \(\rho(X) = \rho(Y)+1\), for example, \(\bbP^n\)-bundles over \(Y\);
            \item \(Y\) is an abelian variety or of picard number one;
            \only<6->{\alert{\item \(f\) has a Zariski dense orbit;}}
            \only<6->{\alert{\item \(\delta_f > \delta_g\);}}
            \only<6->{\alert{\item if \(Y\) is of picard number one, then \(\delta_g = 1\).}}
        \end{itemize}
    }

    \only<2-5>{
        \begin{theorem}[{[Meng-Wang-Y]}]\label{thm: mainthm}
            KSC holds for \((X,f)\) under the above settings.
        \end{theorem}
    }

    \only<3>{
        \vspace{0.5em}
        \alert{
            Essentially generalize the following known results:
            \begin{itemize}
                \item {[Li-Matsuzawa 2021, Theorem 4.1]}, projective bundles on smooth Fano varieties of Picard number one;
                \item {[Nasserden-Zotine 2021, Theorem 1.4]}, projective bundles on elliptic curves.
            \end{itemize}
        }
    }

    \only<4-5>{
        \vspace{0.5em}
        \alert<4>{
            Sketch of idea:
            If \(\delta_f = \delta_g\), then 
            \vspace{0.5em}
            \[ \text{\shortstack{KSC for abelian varieties\\ or polarized endomorphisms}} \implies \text{KSC for } (Y,g) \implies \text{KSC for } (X,f). \]
        }
        \vspace{0.5em}
        \visible<5>{ \alert<5>{If \(\delta_f > \delta_g\), under some allowable extra conditions, we have another fibration by dynamical Iitaka theory.}}
    }

    \only<6>{
        \begin{theorem}[{[Meng-Wang-Y]}]\label{thm: mainthm2}
            There is equivariant fibrations
            \[ f \acts X \overset{\varphi_{f,R_f}}{---\!\twoheadrightarrow} Z \racts h \]
            with \(\dim Z>0\) and \(h\) is \(\delta_f\)-polarized and hence \(\delta_h = \delta_f\).
        \end{theorem}
    }
\end{frame}


\begin{frame}{Dynamical Iitaka Theory}

    \vspace{0.7em}

    \begin{theorem}[{[Meng-Zhang]}]
        \vspace{1em}
        Suppose that \(f^*R_f \equiv q R_f\) for some \(q > 1\).
        Then there exists an \(f\)-equivariant fibration (dynamical Iitaka fibration associated to \(R_f\))
        \[ f \acts X \overset{\varphi_{f,R_f}}{---\!\twoheadrightarrow} Z \racts h \]
        with \(\dim Z = \kappa_f(X,R_f)\) and \(h\) is \(q\)-polarized.
    \end{theorem}

    \vspace*{1em}
    \pause
    \alert<2>{
        With help of this result, we only need to show that under our settings,
        \begin{itemize}
            \item \(\kappa_f(X,R_f) \geq 1\); \visible<3->{ \alert<3>{\(\Longleftarrow\) criterion for toric bundles}}
            \item \(f^*R_f \equiv \delta_f R_f\). \visible<4>{ \alert<4>{\(\Longleftarrow\) decomposition of cones}}
        \end{itemize}
    }

    % \pause
    % \alert<3>{
    %     Then 
    %     \[ \text{KSC for polarized endomorphisms} \implies \text{KSC for } (Y,g) \implies \text{KSC for } (X,f). \]
    % }

\end{frame}


\begin{frame}[c]{Dynamical Iitaka Theory}

    \begin{columns}[c]  
        \begin{column}{0.4\textwidth}
            Coarse classification of varieties via Kodaira dimension \(\kappa(X,K_X)\):
        \end{column}

        \begin{column}{0.6\textwidth}
            \centering
            \begin{tabular}{c|c}
                \hline
                $\kappa(X,K_X)$         & Typical geometry  \\
                \hline
                $-\infty$               & uniruled          \\
                $0$                     & Calabi-Yau        \\
                $0<\kappa(X)<\dim X$    & fibrations type   \\
                $\dim X$                & general type      \\
                \hline
            \end{tabular}
        \end{column}
    \end{columns}

    \pause
    \vspace{0.6cm}

    \begin{columns}[c]
        \begin{column}{0.4\textwidth}
            Dynamical analogue: classify \((X,f)\) via \(\kappa_f(X,R_f)\).
        \end{column}

        \begin{column}{0.6\textwidth}
            \centering
            \begin{tabular}{c|c}
                \hline
                $\kappa_f(X,R_f)$                     & Typical dynamics                      \\
                \hline
                $0$                                   & \(f\) is log-\'etale                  \\
                \alert{$0<\kappa_f(X,R_f)<\dim X$}    & \alert{fibrations type}                       \\
                $\dim X$                              & \(R_{f^s}\) is big for \(s \gg 0\)   \\
                \hline
            \end{tabular}
        \end{column}
    \end{columns}

\end{frame}


\begin{frame}{Technique: decomposition of cones}

    \vspace{0.5em}

    \begin{theorem}[{[Meng-Wang-Y]}]\label{thm: cone-decomp}
        \vspace{0.5em}
        \begin{columns}[c]
            \begin{column}{0.4\textwidth}
                \centering
                \(\pi:(X,f) \to (Y,g)\) flat; \\
                \(X,Y\) smooth projective; \\
                \(\delta_f > \delta_g\);\(\rho(X) = \rho(Y)+1\).
            \end{column}

            \begin{column}{0.1\textwidth}
                \centering
                \[ \Longrightarrow \]
            \end{column}

            \begin{column}{0.5\textwidth}
                \centering
                \(\exists D\) nef with \(f^*D \equiv \delta_f D\) s.t.
                \[ \Psef^1(X) = \pi^*\Psef^1(Y) \oplus \bbR_{\geq 0} D,\] 
                \[ \Nef^1(X) = \pi^*\Nef^1(Y) \oplus \bbR_{\geq 0} D. \]
            \end{column}
        \end{columns}
    \end{theorem}

    \pause

    \begin{columns}[c]  
        \begin{column}{0.5\textwidth}
            \centering

            Cones of a general fibration \(\pi:X\to Y\):

            \vspace{1em}
            
            \begin{tikzpicture}[scale=1, every node/.style={font=\small}]
                % parameters
                \def\rrrr{2.2}
                \def\dddd{0.5}

                % semicircle (Psef^1(X))
                \fill[cyan!8,draw=cyan,thick] (-\rrrr,0) arc (180:0:\rrrr) -- cycle;
                % quadrangle (Nef^1(X))
                \fill[magenta, fill opacity=0.3, draw=magenta, thick] (-\rrrr,-0.03) -- (-0.5*\rrrr,0.866*\rrrr) -- (0,\rrrr) -- (\dddd,-0.03) -- cycle;
                \node[below,magenta!80!black,font=\scriptsize] at (-0.8,0.9) {\(\Nef^1(X)\)};
                % overall label for the semicircle
                \node[above,cyan!70!black] at (0.3,0.7) {\(\Psef^1(X)\)};

                % diameter line segment (Psef^1(Y)) (Nef^1(Y))
                \draw[thick,blue] (-\rrrr,0) -- (\rrrr,0);
                \node[below,blue!80!black] at (0.7,-0.15) {\(\pi^*\Psef^1(Y)\)};
                \draw[thick,red] (-\rrrr,-0.03) -- (\dddd,-0.03);
                \node[below,red!80!black,font=\scriptsize] at (-1.1,0) {\(\pi^*\Nef^1(Y)\)};
                
            \end{tikzpicture}
        \end{column}

        \pause

        \begin{column}{0.5\textwidth}
            \centering

            Cones of fibration with dynamical restrictions:

            \vspace{0.5em}

            \begin{tikzpicture}[scale=1, every node/.style={font=\small}]
                % parameters
                \def\rrrr{2.2}
                \def\dddd{0.5}

                % triangle (Psef^1(X))
                \fill[cyan!8,draw=cyan,thick] (-\rrrr,0) -- (0,\rrrr) -- (\rrrr,0) -- cycle;
                % triangle (Nef^1(X))
                \fill[magenta, fill opacity=0.3, draw=magenta, thick] (-\rrrr,-0.03) -- (0,\rrrr) -- (\dddd,-0.03) -- cycle;
                \node[below,magenta!80!black,font=\scriptsize] at (-0.8,0.9) {\(\Nef^1(X)\)};
                % overall label for the semicircle
                \node[above,cyan!70!black] at (0.3,0.7) {\(\Psef^1(X)\)};

                % diameter line segment (Psef^1(Y)) (Nef^1(Y))
                \draw[thick,blue] (-\rrrr,0) -- (\rrrr,0);
                \node[below,blue!80!black] at (0.7,-0.15) {\(\pi^*\Psef^1(Y)\)};
                \draw[thick,red] (-\rrrr,-0.03) -- (\dddd,-0.03);
                \node[below,red!80!black,font=\scriptsize] at (-1.1,0) {\(\pi^*\Nef^1(Y)\)};

                % dot and label at (0,\rrrr)
                \fill[green] (0,\rrrr) circle (0.06);
                \node[above right,font=\scriptsize,green!70!black] at (0,\rrrr) {\(D\) };
            \end{tikzpicture}
        \end{column}
    \end{columns}
\end{frame}


\begin{frame}{Technique: criterion for toric bundles}
    \vspace{0.5em}
    \begin{theorem}[{[Meng-Wang-Y]}]\label{thm: toric-bundle-criterion}
        % \vspace{0.5em}
        Let \(\pi:(X,f) \to (Y,g)\) be a Fano fibration with \(X\) smooth.
        Suppose that \(\exists\) reduced divisor \(D\) on \(X\) 
        with \alert{\(f^{-1}(D) = D,f^*D \sim q D\)} for some \(q > 1\) and \alert{\(K_X + D \equiv_\pi 0\)}.
        % Then \(\pi:X \to Y\) is a \alert{toric bundle.}
        Then \((X_y,D_y)\) is a \alert{toric pair} for general \(y \in Y\).
    \end{theorem}

    \pause
    \vspace{1em}

    \begin{columns}[t]
        \begin{column}{0.55\textwidth}\vspace{0pt}
            \alert<2>{
                This theorem generalizes the following absolute case (\(Y = \{\mathrm{pt}\}\)) criterion:
                \begin{itemize}
                    \item {[Hwang-Nakayama]}  \(\rho(X) = 1\) case;
                    \item {[Meng-Zhang]}   polarized case;
                    \item {[Meng-Zhong]}   int-amplified case.
                \end{itemize}
            }
            \only<3>{
                \alert<3>{The right sides is an example in the absolute case.}
            }
        \end{column}

        \pause

        \begin{column}{0.4\textwidth}\vspace{0pt}
            \centering
            \begin{tikzpicture}[scale=0.7]
                % Define equilateral triangle vertices
                \coordinate (A) at (0,0);
                \coordinate (B) at (4,0);
                \coordinate (C) at (2,{2*sqrt(3)});
                
                % Calculate extended points for lines beyond vertices
                \coordinate (A') at ($(A)!1.3!(B)$);
                \coordinate (A'') at ($(B)!1.3!(A)$);
                
                \coordinate (B') at ($(B)!1.3!(C)$);
                \coordinate (B'') at ($(C)!1.3!(B)$);
                
                \coordinate (C') at ($(C)!1.3!(A)$);
                \coordinate (C'') at ($(A)!1.3!(C)$);
                
                % Draw extended sides of triangle
                \draw[thick,red] (A'') -- (A') node[below,midway,text=gray!50,font=\scriptsize] {\(Y=0\)};
                \draw[thick,red] (B'') -- (B') node[right,midway,text=gray!50,font=\scriptsize] {\(X=0\)};
                \draw[thick,red] (C'') -- (C') node[left,midway,text=gray!50,font=\scriptsize] {\(Z=0\)};
                
                % Blue arrows along boundary (D is totally invariant)
                \draw[->,thick,blue!60,>=Stealth] ($(A)!0.3!(B)$) -- ($(A)!0.45!(B)$);
                \draw[->,thick,blue!60,>=Stealth] ($(A)!0.7!(B)$) -- ($(A)!0.85!(B)$);
                \draw[->,thick,blue!60,>=Stealth] ($(B)!0.25!(C)$) -- ($(B)!0.4!(C)$);
                \draw[->,thick,blue!60,>=Stealth] ($(B)!0.65!(C)$) -- ($(B)!0.8!(C)$);
                \draw[->,thick,blue!60,>=Stealth] ($(C)!0.3!(A)$) -- ($(C)!0.45!(A)$);
                \draw[->,thick,blue!60,>=Stealth] ($(C)!0.7!(A)$) -- ($(C)!0.85!(A)$);
                
                % Cyan arrows parallel to blue (interior flow)
                \draw[->,thick,cyan!70,>=Stealth] ($($(A)!0.3!(B)$)+(0,0.4)$) -- ($($(A)!0.45!(B)$)+(0,0.4)$);
                \draw[->,thick,cyan!70,>=Stealth] ($($(A)!0.7!(B)$)+(0,0.4)$) -- ($($(A)!0.85!(B)$)+(0,0.4)$);
                \draw[->,thick,cyan!70,>=Stealth] ($($(B)!0.25!(C)$)+(-0.35,-0.2)$) -- ($($(B)!0.4!(C)$)+(-0.35,-0.2)$);
                \draw[->,thick,cyan!70,>=Stealth] ($($(B)!0.65!(C)$)+(-0.35,-0.2)$) -- ($($(B)!0.8!(C)$)+(-0.35,-0.2)$);
                \draw[->,thick,cyan!70,>=Stealth] ($($(C)!0.3!(A)$)+(0.35,-0.2)$) -- ($($(C)!0.45!(A)$)+(0.35,-0.2)$);
                \draw[->,thick,cyan!70,>=Stealth] ($($(C)!0.7!(A)$)+(0.35,-0.2)$) -- ($($(C)!0.85!(A)$)+(0.35,-0.2)$);
                
                % Cyan arrows 
                \draw[->,thick,cyan!70,>=Stealth] ($($(A)!0.3!(B)$)+(0.8,0.6)$) -- ($($(A)!0.45!(B)$)+(0.8,0.6)$);
                \draw[->,thick,cyan!70,>=Stealth] ($($(B)!0.25!(C)$)+(-1.05,0.4)$) -- ($($(B)!0.4!(C)$)+(-1.05,0.4)$);
                \draw[->,thick,cyan!70,>=Stealth] ($($(C)!0.3!(A)$)+(0.15,-1.05)$) -- ($($(C)!0.45!(A)$)+(0.15,-1.05)$);

                % Mark vertices with dots
                \fill[purple] (A) circle (0.05);
                \fill[purple] (B) circle (0.05);
                \fill[purple] (C) circle (0.05);
                
                % Label vertices
                \node[below left,gray!50,font=\scriptsize] at (A) {\([1:0:0]\)};
                \node[below right,gray!50,font=\scriptsize] at (B) {\([0:1:0]\)};
                \node[above,gray!50,font=\scriptsize] at (C) {\([0:0:1]\)};
            \end{tikzpicture}

            \vspace{-0.5em}

            {\scriptsize
            \begin{align*}
                f: & \bbP^2 \to \bbP^2, \\
                   & [X:Y:Z] \mapsto [X^q:Y^q:Z^q]; \\
                D= & \{XYZ=0\}.
            \end{align*}
            }
        \end{column}
    \end{columns}
    
\end{frame}


\begin{frame}[c]{Further questions}

    Here are two further questions we are interested in:

    \begin{itemize}
        \visible<2->{\alert<2>{\item weaken the conditions on \(Y\): can we just assume \(Y\) is normal projective?}}
        \visible<3->{\alert<3>{\item deal the case \(R_f\) is big in the dynamical Iitaka theory.}}
    \end{itemize}
    
\end{frame}


% ----------------------------------------------------------

% Thank you slide
{
    \setbeamertemplate{background}{}
    \begin{frame}[plain]

    \vspace{0.2\textheight}
    \Huge{\centerline{\textbf{\alert{Thank You!}}}}

    \vfill
    \centerline{\includegraphics[width=0.35\textwidth]{\ECNULibThanks}}
    \end{frame}
}


\end{document}